\chapter{BioFeaturisers}
\label{ch:biofeaturisers}

% In this chapter we introduce the BioFeaturisers library, which is a collection of differentiable biophysical forward models. These models can be used with jaxENT but also include pytorch wrappers. This way we have a unifed approach to forward models that can be used in both metadynamics as well as metadiffusion. This is made possible by a heavy-side contact BV-model formulation that is differentiable as well an optimised SAXS forward model using FFT. First case study we use BPTI as a test case with the 1ms Anton trajectory for metainference-metadynamics simulations. We then apply to this to MoPrP under pH 4 and pH 7, we use M&M to generate initial candidate ensemble and perform unrestrained short-APO simulations from decorrelated structures to observe pH-dependant unfolding pathways. 
% The second part applies the pipeline to metadiffusion, where we show how we can combine diffusion-based models with experimental guidance to perform ensemble generation of complexes. Second case study looks at ligand-induced dynamical changes with RoRgamma that cover the full range of agonist, partial-agonist, and antagonist ligands. We compare the results of metadiffusion with metadynamics and show that we can use the same forward models to perform both ensemble generation and ensemble reweighting. 

% Define heavy-side BV-contacts, Define SAXS forward model + FFT version, Define metainference, metadynamics and then metadiffusion.

\section{Motivation}
\label{sec:biofeat-motivation}

\subsection{Abstract}
\label{sec:biofeat-abstract}

\section{Introduction}
\label{sec:biofeat-intro}

\section{Methods}
\label{sec:biofeat-methods}

\subsection{Heaviside Contact BV Model}
\label{sec:biofeat-bv}

\subsection{SAXS Forward Model}
\label{sec:biofeat-saxs}

\subsubsection{FFT SAXS}
\label{sec:biofeat-saxs-fft}

\subsection{Metainference}
\label{sec:biofeat-metainference}

\subsection{Metadynamics}
\label{sec:biofeat-metadynamics}

\subsection{Metadiffusion}
\label{sec:biofeat-metadiffusion}
